\section{Justification beyond proof}\label{sec:philosophy}

The investigation began with a question about principled justification. What does it look like, once we go beyond mathematical proof? This section steps back from the theorems and says what they show about that question. It also says what they do not show. Throughout, ``we'' means the authors of this report. The view argued for is deliberately modest. The theorems do not deliver a theory of justification. They make precise one shape that justification can take, and they locate exactly where that shape runs out.

\subsection{The setup is Goodman's circle, made into a learning problem}\label{sec:philosophy:goodman}

The setup we studied has three parts:
\begin{itemize}
\item learn inference rules from accepted inferences;
\item revise them when they generate unacceptable conclusions;
\item revise accepted inferences when they conflict with rules one is unwilling to give up.
\end{itemize}
This is, almost verbatim, Goodman's account of the justification of deduction and induction:
\begin{quote}
``Rules and particular inferences alike are justified by being brought into agreement with each other. \emph{A rule is amended if it yields an inference we are unwilling to accept; an inference is rejected if it violates a rule we are unwilling to amend.}'' \citep[ch.~III]{goodman1955fact}
\end{quote}
In our vocabulary, the channels correspond to Goodman's account as follows:
\begin{itemize}
\item Imitation (P) supplies the accepted particular inferences.
\item Rule schemas, learned by least generalization, are the rules.
\item Coherence (C) is the mechanism by which a rule ``yields an inference we are unwilling to accept''. A derivation of $\bot$ from a context certified consistent is a negative bag, and some rule in it must go.
\item World feedback (W) adds a channel whose errors are independent of human errors.
\end{itemize}
Without (W), the process is \emph{narrow} reflective equilibrium \citep{rawls1971theory,rawls1974independence}. With (W) it is \emph{wide} reflective equilibrium in the sense of \citet{daniels1979wide}. Daniels' ``independence constraint'' requires that the background theories not rest on the same considered judgments as the principles. It can be read as requiring that the corroborating channel's errors be independent of the human ones. Without such independence, agreement moves the odds by at most a bounded factor (\cref{sec:coherence:eliminative}). Even with it, no informative coherence measure is truth-conducive in general.

What our theorems add to Goodman's picture is a set of \emph{vindications with explicit conditions}. Under stated channel assumptions, mutual adjustment is \emph{sound against every adaptive arguer}, with explicit costs:
\begin{itemize}
\item realizability of the target in a structured hypothesis class;
\item truthful designation of the contexts that are certified consistent;
\item a world channel of the right kind.
\end{itemize}
For several of these assumptions there is a matching impossibility result (\cref{tab:twotier:necessary}). Realizability is shown to be load-bearing by example rather than by a theorem. The world channel is needed for completeness under ambiguous blame, not for soundness (\cref{thm:twotier:blame,prop:twotier:coherenceonly}). This is the sense in which the theorems are \emph{conditional reliabilist meta-justifications} of reflective equilibrium. They are not a non-circular foundation for logic. As \citet{dummett1973justification} observed, any justification of deduction is explanatory rather than suasive: it uses deduction. What they provide is an explicit list of what the stated guarantees rely on, with proofs that the guarantees hold under them. The list consists of the three items above, together with the technical hypotheses of \cref{thm:twotier:main}: a known frequency floor, fallacies used under their own latent tags, and an exact depth-$d$ refutation oracle.

Two features of the guarantees matter philosophically.
\begin{itemize}
\item First, they are \emph{anytime}, not merely in the limit. Limit theorems are Reichenbachian vindications \citep{reichenbach1938experience}. They fall to the objection of \citet{salmon1967foundations} that infinitely many methods share the same limit while disagreeing at every finite stage. Our soundness guarantees hold at every time, with probability over the data (\cref{thm:caution:ville}). Our costs are explicit finite quantities: escalations, detections, samples. The motivating notes are suspicious of limit framings, and hold that thinking and mathematics are ``infinite endeavors'' with no final formula. They also propose ``trajectories of moral-reflective flight'' as an alternative to reflective equilibrium, which they find ``pretty flat/dead'' as a picture of a mind's development (\texttt{philosophy/ethics/Trajectories of moral-reflective flight \textemdash{} an alternative to reflective equilibrium.md}). Our results bear on that view only in a limited way. They concern a \emph{fixed} target calculus, for which convergence is the right notion. They say nothing about the open-ended development of concepts and calculi that the notes have in mind (language invention, \cref{sec:open}).
\item Second, they are \emph{worst case over arguers}. A step checker that is accurate on the distribution of human steps can be exploited by search (\cref{thm:search:tonk}). Justification that is to survive an untrusted or adversarial generator must be uniform over what the generator may propose. This is the case that motivates the question in the first place.
\end{itemize}

\subsection{What each signal can do, and the residue it leaves}\label{sec:philosophy:labour}

The main results combine into a division of labour. \Cref{tab:philosophy:labour} summarizes it.

\begin{table}[ht]
\centering\small
\begin{tabularx}{\textwidth}{>{\raggedright\arraybackslash}p{2.2cm}>{\raggedright\arraybackslash}X>{\raggedright\arraybackslash}X}
\toprule
Signal & What it provably fixes & What it provably cannot fix (residue)\\
\midrule
Imitation (P) & The \emph{rules} actually followed. Exact identification of a schematic calculus at coupon-collector rates (\cref{thm:imitation:coupon}), with systematic out-of-distribution generalization (\cref{cor:imitation:ood}). & Systematic errors, which are indistinguishable from rules at every rate (\cref{cor:imitation:systematic}). Imitation converges to ``competence'' in the sense of \citet{cohen1981can}, not to validity \citep{stich1980justification}.\\
Cautious acceptance & Soundness against every adaptive arguer (\cref{thm:caution:vs,thm:caution:ville}). & Nothing new. It receives no coherence signal at all (\cref{prop:coherence:caution}), and its price is set by the structure of the class (\cref{tab:caution:cost}).\\
Coherence (C) & Every error that is incoherent with the target on a designated context (\cref{thm:coherence:elimination}). Exactly the target, when it has no coherent uniform proper extension: classical logic (\cref{thm:coherence:cpc}), and complete theories at the level of theorems. The classical \emph{meanings} of the connectives, once denials are allowed (\cref{thm:coherence:telltale}). & Coherent uniform alternatives: the Kripkensteinian residue (\cref{thm:coherence:residue}). Intermediate logics (\cref{thm:coherence:ipc}). $\Sigma_1$-unsound extensions of arithmetic (\cref{thm:coherence:popper}). Which valuation is actual.\\
Simplicity & Expensive, idiosyncratic errors, via a rate threshold (\cref{sec:simplicity}). & Cheap, frequent fallacies. Rare valid rules are lost. Its strength cannot be calibrated from imitation data.\\
World (W) & Objects give step-level blame by descent, so $\log_2|\Hc|$ corrections suffice, against up to $|\Hc|-1$ from paradoxes alone on some classes (\cref{thm:informal:objects,thm:informal:paradoxes}). Computation with fresh randomness gives worst-case-valid stochastic checks (\cref{lem:search:commitment}). Calibration of top models (\cref{sec:physics}). & Errors invisible on every presentable object. Effects outside the top model (\cref{thm:physics:judgment}).\\
Contexts & Reasoning under contradictory idealizations without explosion. Coherence where it belongs (\cref{thm:physics:realizability}). & The reading of a problem, and the adequacy of the top model. Both are ineliminable judgments (\cref{thm:physics:judgment}).\\
\bottomrule
\end{tabularx}
\caption{The division of labour among the learning signals, with the residue each one provably leaves.}\label{tab:philosophy:labour}
\end{table}

Two observations about this table seem to us to be the real philosophical content of the investigation.

\paragraph{Every signal is one-sided.}
Each learning signal we studied can refute errors in only one direction:
\begin{itemize}
\item Positive data never refute an over-generalization (Gold).
\item Coherence never refutes over-caution: a cautious learner never sees a contradiction (\cref{prop:coherence:caution}). Nor does it refute a coherent over-generalization (\cref{thm:coherence:silent}).
\item Coherence among exported tolerances never refutes a tolerance that is too wide. Widening an interval cannot create a contradiction (\cref{thm:physics:coherencetol}).
\item The demand that a problem be well posed refutes readings that are too wide, but it rewards readings that are too narrow (\cref{prop:physics:gricean}).
\end{itemize}
Each one-sided signal is an attack surface for an adversarial arguer, and it calls for a complementary signal that refutes errors in the other direction. Principled justification, on this picture, is not one master criterion. It is a set of signals that \emph{together} cover every direction of error, plus an explicit statement of the directions that remain uncovered. \Cref{tab:physics:onesided} lists the pairs for physics. The same structure appears in mathematics. Imitation proposes and caution guards. Coherence condemns, but only bold hypotheses. Objects localize.

\paragraph{Each residue has a name, and different philosophies disagree only about the residue.}
Outside the residue of coherent, uniform, world-adequate alternatives, every position we know of counts the learner as correct:
\begin{itemize}
\item the realist;
\item the Peircean, for whom truth is the limit of inquiry;
\item the Dummettian, for whom truth is warranted assertibility;
\item Carnap's tolerance;
\item Quine's pragmatism.
\end{itemize}
Inside the residue, they disagree. The realist sees undetectable error. The Peircean or Dummettian sees no fact of the matter. Carnap sees a free choice of language. Quine sees a choice made on pragmatic grounds. The theorems make the residue precise in several cases:
\begin{itemize}
\item It is empty for classical propositional logic.
\item It contains all intermediate logics for an intuitionistic practice.
\item It contains $\PA+\neg\Con(\PA)$ for arithmetic.
\item In the algebra experiment it contains the total semantics with $1/0=0$, which pure coherence converged to (\cref{sec:experiments}). This is, as it happens, exactly the convention adopted by Lean's Mathlib and by Isabelle/HOL, where $x/0=0$.
\end{itemize}
A theorem can say exactly where these philosophies part ways. It cannot settle who is right inside the residue. We think that is the correct division of labour between mathematics and philosophy here.

One consequence deserves to be stated on its own. \emph{Rules can be identifiable when truths are not.} The bilateral tell-tale identifies classical meanings from finitely many data (\cref{thm:coherence:telltale}). Yet no computable learner on the same channels decides $\Sigma_2$ truth in the limit (\cref{thm:coherence:popper}). This bears on Dummett's acquisition argument but does not refute it. That argument targets truth-conditional meaning: a meaning one could acquire must be manifestable in use, so meaning cannot outrun decidability. What the learner acquires is the rules, for the propositional connectives (\cref{thm:coherence:telltale}), and Dummett agrees that rules are manifestable. The theorems show only that acquiring rules exactly is compatible with truth values, here of arithmetic $\Sigma_2$ sentences, staying beyond computable reach.

\subsection{Rules are not premises}\label{sec:philosophy:carroll}

The tortoise of \citet{carroll1895tortoise} refused to infer $Z$ from $A$, $B$ and ``if $A$ and $B$ then $Z$'' until that inference was itself added as a premise, and so on forever. For a learner the lesson has a precise form:
\begin{itemize}
\item A learned \emph{rule} is a schema. Its validity is a universally quantified claim over all instances, of $\Pi$ form.
\item No finite amount of world feedback verifies such a claim. Feedback can only refute it.
\item The justification for \emph{using} a learned rule therefore cannot be further premises. It has to be a guarantee about the process that accepted it: cautious acceptance, and an anchor that has been observed.
\end{itemize}
This is why the right acceptance policy for learned rules is Popperian. Accept tentatively. Assert only what every surviving hypothesis licenses. Treat each refutation as permanent. It is also why self-certification must never feed back into acceptance:
\begin{itemize}
\item a learned verifier that ``proves its own soundness'' has provided no evidence that distinguishes it from an unsound verifier sharing the certifying rule;
\item in arithmetic, by Löb's theorem a theory that proves the reflection schema $\mathrm{Prov}_T(\ulcorner\varphi\urcorner)\to\varphi$ for its \emph{own} provability predicate, for all $\varphi$, proves everything. By contrast, the corresponding reflection \emph{rule} is admissible for $\Sigma_1$-sound theories.
\end{itemize}
Reflection is a principled source of genuinely new inferences, beyond imitation, that a learner who has accepted a calculus is committed to. Where it is used, it must enter as a stratified rule, never as a self-referential premise.

\subsection{What principled justification beyond proof looks like}\label{sec:philosophy:shape}

Putting the pieces together, the results suggest the following shape for a principled justification of a claim about something other than pure mathematics. A physics olympiad answer is the running case.

\begin{enumerate}
\item \textbf{Proof inside contexts.} Each step inside a context is checked against a calculus. The calculus is formal where possible. Where it is learned, acceptance is cautious, so the check is sound against an arguer who searches for weaknesses (\cref{sec:caution,sec:informal}).
\item \textbf{Certified bridges between contexts.} Moving a conclusion from an idealized context to the claim of interest is not a logical step. It needs a certificate carrying information of finite radius about how the quantity behaves away from the idealization. Facts about the idealized model alone never suffice (no free export, \cref{thm:physics:nofree}). Side conditions are evaluated where they are true, in the parent context, and contexts are declared deformations. Otherwise an arguer can ``stipulate'' its way to anything (\cref{thm:physics:stipulation}).
\item \textbf{An explicit judgment layer.} What remains is provably ineliminable (\cref{thm:physics:judgment}):
  \begin{itemize}
  \item that the problem was read correctly, i.e.\ the admissible worlds were not narrowed;
  \item that the top model is adequate for the world claims;
  \item that the trusted base is trustworthy.
  \end{itemize}
  A principled justification does not hide this layer. It minimizes it and declares it.
\item \textbf{Correctness as supertruth.} For underspecified questions, ``true in the context at hand'' means true in every admissible completion of the intended model, to the demanded precision. This is the largest acceptance criterion that stays sound under uncertainty about the reading (\cref{thm:physics:superval}). Answer-invariance under unspecified features, which experienced solvers notice and rely on, is exactly a certificate of supertruth.
\item \textbf{Learning signals with known limits.} Every learned component comes with a statement of the signal that trained it and the residue that signal leaves. The learned components are the calculus, the bridge validity regions and the tolerances.
\end{enumerate}

Measured against this shape, mathematical proof is the limiting case in which:
\begin{itemize}
\item the judgment layer has shrunk to the reading of the formal statement and trust in the axioms;
\item every bridge is exact;
\item the calculus is fixed rather than learned.
\end{itemize}
Informal mathematics before formalization sits in between. In our model it is a bounded-gap practice over a latent calculus. If the robust-core hypothesis holds, its proofs survive every admissible formalization. In the historical episodes we examined, counterexample objects did most of the corrective work, and the few paradoxes were the set-theoretic antinomies (\cref{sec:informal}). Physics adds a nonempty, explicit judgment layer and non-logical bridges, but the shape is the same. On this view the difference between mathematical and physical justification is one of degree, measured by the size of the declared residue, not one of kind.

\subsection{What this says about meaning}\label{sec:philosophy:meaning}

The inferentialist premise of the question was that learning which inferences are valid is roughly learning meanings. The results support a qualified version:
\begin{itemize}
\item \textbf{Inferential roles are learnable, exactly, from finite use.} A finite anchor fixes the inferential role of a vocabulary for unbounded use (\cref{cor:imitation:ood}). This is a precise, if modest, answer to the rule-following question of how finite training fixes an infinite norm. The answer is: through structure (schemas) plus a small set of telling instances. The skeptic's ``quus'' survives only as a non-uniform patch, which uniformity rules out, or as a coherent uniform alternative, which is a named residue. This answer is relative to the hypothesis class and its vocabulary, which the data do not fix. That is Goodman's point, and it is exactly where Kripke's skeptic presses (\cref{sec:imitation:ood}).
\item \textbf{More than role is not identifiable.} Under complete closure-level data, practice, coherence and object data identify a latent formalization exactly up to informal-validity equivalence, and no further. Bounded-gap data only bracket it (\cref{thm:informal:identify}). Ontology is not identifiable: standard and nonstandard reals are observationally equivalent on the practice (\cref{prop:informal:transfer}). Neither is the choice of set-theoretic reduction (\cref{prop:informal:benacerraf}).
\item \textbf{Coherence buys existence, but not intended existence.} Every coherent mode of talking has \emph{something} it could be talking about. That is a generalized completeness theorem, true for positions, credences and context systems alike (\cref{sec:existence}). But that something may be a syntactic or nonstandard model. Intended existence needs one of three additional things:
  \begin{itemize}
  \item categoricity;
  \item anchoring of the vocabulary to the world (Beth definability);
  \item an infinitary condition, such as the probabilistic $\omega$-rule.
  \end{itemize}
  None of these can be certified computably beyond $\Pi_1$.
\end{itemize}
Inferential role is one aspect of meaning, and the project captures that aspect precisely. Other aspects are not captured by anything studied here:
\begin{itemize}
\item the ``hooking'' of a conception onto the world;
\item the activities a concept supports;
\item the questions it lets one ask.
\end{itemize}
We do not claim otherwise. All three aspects are emphasized in the motivating notes. These treat inferential role (``compositional inference warranting'') as one aspect of meaning among several (\texttt{philosophy/philosophy of language/meaning and truth/}; \texttt{philosophy/philosophy of thinking/beating solomonoff induction at grokking a notion.md}).

\subsection{The motivating questions}\label{sec:philosophy:answers}

The questions that prompted this investigation are recorded in H\"anni's public working notes \citep{hanni2026notes}. \Cref{sec:intro:answers} and \cref{tab:intro:answers} list them, with pointers to where each is answered or partly answered. The readings of the notes given there are ours.
